% ============================================================
% Thesis Chapter: Energy Profiling of 5G User Plane Functions
% ============================================================
% Compile with:  pdflatex chapter_upf_profiling.tex  (run twice)
% ============================================================

\documentclass[12pt, a4paper]{report}

% ── Packages ────────────────────────────────────────────────────────────────
\usepackage[T1]{fontenc}
\usepackage[utf8]{inputenc}
\usepackage{lmodern}
\usepackage{amsmath, amssymb}
\usepackage{booktabs}
\usepackage{graphicx}
\usepackage{subcaption}
\usepackage{xcolor}
\usepackage[margin=2.5cm]{geometry}
\usepackage[hidelinks]{hyperref}
\usepackage{setspace}
\onehalfspacing

% ── Convenience macros ──────────────────────────────────────────────────────
\newcommand{\dpdk}{\textsc{sd-core dpdk}}
\newcommand{\usr}{\textsc{usr-upf}}
\newcommand{\upf}{\textsc{upf}}
\newcommand{\secm}{$\mathrm{SEC}$}
\newcommand{\sectable}{$\mathrm{SEC}_{\mathrm{total}}$}
\newcommand{\secnet}{$\mathrm{SEC}_{\mathrm{net}}$}
\newcommand{\figpath}[1]{figures/#1}

\begin{document}
% ============================================================

\chapter{Energy Profiling of 5G User Plane Functions}
\label{chap:upf_profiling}

% ── Section 1: Introduction ─────────────────────────────────────────────────
\section{Introduction}
\label{sec:intro}

The User Plane Function (\upf{}) is the component of a 5G core network
responsible for routing and forwarding GTP-U encapsulated user traffic between
the Radio Access Network (RAN) and the Data Network (DN).
Its power consumption is therefore directly coupled to network load, and
understanding this relationship is a prerequisite for energy-aware network
management and for meeting increasingly stringent sustainability targets in
mobile telecommunications.

This chapter presents a systematic empirical study of two open-source \upf{}
implementations under controlled laboratory conditions.
The first implementation, \dpdk{}, is based on the SD-Core project and
leverages the Data Plane Development Kit (DPDK) poll-mode I/O driver.
The second, \usr{}, is the user-plane component of OpenAirInterface and
uses a conventional interrupt-driven Linux I/O model.
Traffic was generated with a Keysight LoadCore emulator, and power was
measured at the host level using the Scaphandre process-level power monitor
with 1\,s sampling resolution.

The remainder of this chapter is organised as follows.
Section~\ref{sec:methodology} describes the measurement methodology and
dataset.
Section~\ref{sec:profiles} characterises each variant across performance
and energy dimensions.
Section~\ref{sec:saturation} investigates the saturation behaviour of
\usr{} in depth.
Section~\ref{sec:sec} analyses the Specific Energy Consumption metric and
its interpretation for poll-mode versus interrupt-driven drivers.
Section~\ref{sec:correlation} reports feature correlations relevant to
predictive energy modelling.
Section~\ref{sec:implications} discusses implications for network operation
and modelling.

% ── Section 2: Methodology ──────────────────────────────────────────────────
\section{Measurement Methodology}
\label{sec:methodology}

\subsection{Testbed Configuration}

Both \upf{} variants were deployed on identical server hardware.
LoadCore injected downlink UDP flows with a fixed packet size of
1250\,B at progressively increasing offered loads.
Each load step was sustained for a fixed duration to allow steady-state
measurements; the raw time-series data were then merged with per-second power
readings from Scaphandre.

\subsection{Dataset}

After preprocessing, the merged dataset comprises 9556 one-second
measurement samples across 110 USR runs and an equivalent number of
DPDK runs.
Following feature engineering (removal of zero-variance and near-duplicate
columns, unit conversions, target computation), the analysis dataset contains
8781 samples and 44 columns, split evenly between the two variants
(4388 DPDK, 4393 USR).

\subsection{Metrics}

The following primary metrics are reported throughout this chapter:

\begin{description}
  \item[Power~(W)] Total host power consumption in watts, as reported by
    Scaphandre (\texttt{power\_watts}).
  \item[Net power~(W)] Incremental power above the idle baseline:
    $P_{\mathrm{net}} = P_{\mathrm{total}} - P_{\mathrm{idle}}$.
  \item[\sectable{}] Specific Energy Consumption based on total power:
    $\mathrm{SEC}_{\mathrm{total}} = P_{\mathrm{total}} / \lambda$,
    where $\lambda$ is throughput in Gbps.
  \item[\secnet{}] Marginal energy cost: $\mathrm{SEC}_{\mathrm{net}}
    = P_{\mathrm{net}} / \lambda$.
  \item[Packet loss] Count of GTP-U packets lost per one-second interval
    at the DN-side receiver.
  \item[Downlink delay] Weighted mean one-way packet delay on the
    downlink path ($\mu$s).
\end{description}

% ── Section 3: Variant Profiles ─────────────────────────────────────────────
\section{Variant Profiles}
\label{sec:profiles}

\subsection{Overview}

Table~\ref{tab:profiles} summarises the key performance and energy
characteristics of both variants across all measurement conditions.
Figure~\ref{fig:power_vs_throughput} illustrates the fundamental
difference in the power--throughput relationship.

\begin{table}[htbp]
  \centering
  \caption{UPF variant profiles: key performance and energy metrics.}
  \label{tab:profiles}
  \setlength{\tabcolsep}{7pt}
  \begin{tabular}{lrr}
    \toprule
    \textbf{Metric} & \textbf{SD-Core DPDK} & \textbf{USR-UPF} \\
    \midrule
    \multicolumn{3}{l}{\textit{Power}} \\
    \quad Total power range (W)         & $0.807$--$0.859$ & $0.000$--$4.497$ \\
    \quad Mean power (W)                & $0.821$           & $1.435$           \\
    \quad Standard deviation (W)        & $0.005$           & $1.527$           \\
    \quad Dynamic range (W)             & $0.052$           & $4.497$           \\
    \quad Net power mean (W)            & $0.001$           & $1.419$           \\
    \midrule
    \multicolumn{3}{l}{\textit{CPU}} \\
    \quad Mean CPU (\%)                 & $1.064$           & $1.549$           \\
    \quad Max CPU (\%)                  & $1.079$           & $4.554$           \\
    \quad Idle CPU (\%)                 & $1.064$           & $0.018$           \\
    \midrule
    \multicolumn{3}{l}{\textit{Throughput}} \\
    \quad Maximum (Gbps)                & $5.06$            & $2.33$            \\
    \quad Mean (Gbps)                   & $0.788$           & $0.528$           \\
    \midrule
    \multicolumn{3}{l}{\textit{Reliability}} \\
    \quad Packet loss (sample rows, \%) & $0.0$             & $32.7$            \\
    \midrule
    \multicolumn{3}{l}{\textit{Latency}} \\
    \quad DL delay mean ($\mu$s)  & $311$      & $5031$            \\
    \quad DL delay max ($\mu$s)   & $24978$    & $25000$           \\
    \quad RAN jitter mean ($\mu$s) & $1.3$     & $1428$            \\
    \midrule
    \multicolumn{3}{l}{\textit{Energy efficiency at $>$1\,Gbps}} \\
    \quad \sectable{} median (J/Gb)     & $0.27$            & $1.61$            \\
    \quad \secnet{} median (J/Gb)       & $0.002$           & $1.61$            \\
    \midrule
    L2/L3 overhead ratio (mean)         & $2.11$            & $1.86$            \\
    \bottomrule
  \end{tabular}
\end{table}

\begin{figure}[htbp]
  \centering
  \includegraphics[width=\linewidth]{\figpath{fig1_power_vs_throughput}}
  \caption{Power consumption versus throughput for each UPF variant.
    SD-Core DPDK (left) exhibits a near-constant power profile due to
    continuous CPU polling. USR-UPF (right) shows a strong load-proportional
    relationship but saturates above 0.5\,Gbps.}
  \label{fig:power_vs_throughput}
\end{figure}

\subsection{SD-Core DPDK}

The DPDK data-plane driver operates in poll mode: a dedicated CPU thread
continuously spins on the receive queue, eliminating interrupt overhead at
the cost of occupying the core regardless of traffic.
This architectural choice manifests clearly in the measurements.

Power consumption remains almost constant at $0.821 \pm 0.005$\,W across
the entire load range (0--5.06\,Gbps), yielding a dynamic range of only
52\,mW.
CPU utilisation is similarly invariant at $1.064 \pm 0.004$\,\%, and the
Pearson correlation between throughput and CPU is $r = 0.016$, confirming
that the polling loop consumes a fixed share of the core irrespective of
packet rate.

Despite its flat power profile, \dpdk{} achieves the highest throughput
of both variants (5.06\,Gbps) with zero packet loss
across all 4388 measurement samples.
Downlink delay is low (mean 311\,$\mu$s) and RAN jitter
is negligible (mean 1.3\,$\mu$s), making it the preferable
choice for latency-sensitive traffic.

\subsection{USR-UPF}

\usr{} uses the standard Linux socket and eBPF/XDP interrupt path.
Power scales strongly with offered load: $r(\text{throughput, power}) =
0.81$ (compared with $0.46$ for DPDK), and the system transitions from
near-zero power (0.034\,W) at idle to a plateau of approximately
3.4\,W once saturated.
CPU utilisation follows the same trend ($r = 0.76$), rising from 0\,\%
at idle to 4.6\,\% at peak.

However, \usr{} saturates at approximately 0.5\,Gbps
(see Section~\ref{sec:saturation}), beyond which virtually all packets are
dropped.
Delay is consistently higher than DPDK (mean 5\,ms) and
RAN jitter degrades catastrophically above saturation (mean
1428\,$\mu$s, maximum 199\,ms).

Figure~\ref{fig:variant_profiles} illustrates the load-stratified
comparison of both variants across the four key metrics.

\begin{figure}[htbp]
  \centering
  \includegraphics[width=\linewidth]{\figpath{fig2_variant_profiles}}
  \caption{Load-stratified variant profiles. Throughput bins on the
    x-axis (Gbps). (a)~Power consumption; (b)~CPU utilisation;
    (c)~packet loss incidence; (d)~mean downlink delay.
    Error bars in (a) denote $\pm 1$ standard deviation.
    DPDK maintains flat power and zero loss at all loads; USR scales
    with load but saturates above 0.5\,Gbps.}
  \label{fig:variant_profiles}
\end{figure}

% ── Section 4: USR Saturation ────────────────────────────────────────────────
\section{USR-UPF Saturation Analysis}
\label{sec:saturation}

\subsection{Saturation Onset}

To characterise the saturation regime with fine granularity, measurements
were binned into 50\,Mbps throughput intervals.
Figure~\ref{fig:oai_saturation} shows the transition across four dimensions:
loss incidence, loss magnitude, delay, and the power/CPU trajectory.

The first bin where more than 5\,\% of samples exhibit packet loss occurs at
$\approx$\,0.17\,Gbps, indicating the onset of
saturation.
Loss increases steeply, reaching 50\,\% at $\approx$\,0.57\,Gbps,
beyond which essentially all samples report loss.
From this point, \usr{} enters a fully saturated regime: throughput
measurements reflect the offered load injected by LoadCore, not the
effective forwarding rate, and delay jumps to the measurement ceiling
of 25\,ms.

\begin{figure}[htbp]
  \centering
  \includegraphics[width=\linewidth]{\figpath{fig3_oai_saturation}}
  \caption{USR-UPF saturation transition in 50\,Mbps
    throughput bins. The dashed vertical line marks the saturation onset at
    $\approx$\,0.17\,Gbps.
    (a)~Fraction of samples with packet loss;
    (b)~mean number of packets lost per interval (when loss occurs);
    (c)~mean downlink delay;
    (d)~mean power and CPU utilisation (dual axis).}
  \label{fig:oai_saturation}
\end{figure}

\subsection{Power Behaviour at Saturation}

Despite losing packets, \usr{} continues to consume power proportionally
to the offered load up to $\approx$\,0.5\,Gbps, after
which power plateaus at approximately 3.4\,W
(Figure~\ref{fig:oai_saturation}d).
This indicates that the CPU is fully committed to packet processing at
saturation; additional offered load does not increase power but increases
loss instead.

The behaviour is summarised as follows:
\begin{itemize}
  \item Pre-saturation ($<$\,0.3\,Gbps):
    mean power $= 0.338$\,W, zero loss.
  \item Transition (0.3--0.57\,Gbps): power rises sharply, loss begins.
  \item Post-saturation ($>$\,0.5\,Gbps):
    mean power $\approx 3.4$\,W, near-total loss.
\end{itemize}

\subsection{Per-Run Characterisation}

Figure~\ref{fig:oai_per_run} aggregates the loss and power behaviour at the
granularity of individual measurement runs.
Of the 110 USR runs, 61 (55\,\%) completed without any packet
loss; these runs had a maximum offered throughput of at most
0.1\,Gbps.
Every run exceeding 0.06\,Gbps maximum offered load
incurred at least some packet loss, establishing the \emph{practical safe
operating boundary} of \usr{} at approximately 100\,Mbps.

\begin{figure}[htbp]
  \centering
  \includegraphics[width=\linewidth]{\figpath{fig7_oai_per_run_saturation}}
  \caption{USR-UPF per-run saturation summary.
    Each point represents one measurement run.
    (a)~Overall packet loss versus maximum offered throughput, coloured
    by mean power; (b)~mean power versus maximum offered throughput,
    coloured by packet loss fraction.}
  \label{fig:oai_per_run}
\end{figure}

% ── Section 5: SEC Analysis ──────────────────────────────────────────────────
\section{Specific Energy Consumption Analysis}
\label{sec:sec}

\subsection{Definitions}

Specific Energy Consumption (\secm{}) is defined as the ratio of power
consumed to useful throughput delivered.
Two variants are considered:
\begin{align}
  \mathrm{SEC}_{\mathrm{total}} &= \frac{P_{\mathrm{total}}}{\lambda}
    \quad \left[J/Gb\right] \label{eq:sec_total} \\
  \mathrm{SEC}_{\mathrm{net}}   &= \frac{P_{\mathrm{net}}}{\lambda}
    = \frac{P_{\mathrm{total}} - P_{\mathrm{idle}}}{\lambda}
    \quad \left[J/Gb\right] \label{eq:sec_net}
\end{align}

\subsection{The Poll-Mode Measurement Paradox}

Equation~\eqref{eq:sec_net} was introduced in the literature to capture
the \emph{marginal} energy cost of transmitting one additional gigabit,
isolating the load-dependent component from the fixed idle overhead.
While this definition is appropriate for interrupt-driven systems (where
$P_{\mathrm{idle}} \ll P_{\mathrm{total}}$), it produces a misleading
result for poll-mode drivers.

In \dpdk{}, the CPU polls continuously regardless of load, so
$P_{\mathrm{idle}} \approx P_{\mathrm{total}} \approx 0.821$\,W at all
times.
Consequently, $P_{\mathrm{net}} = P_{\mathrm{total}} - P_{\mathrm{idle}}
\approx 0$, and \secnet{} approaches zero as throughput increases,
superficially suggesting near-perfect efficiency.
In reality, \dpdk{} consumes 0.82\,W regardless of whether it
forwards 0 or 5\,Gbps.

Table~\ref{tab:sec} and Figure~\ref{fig:sec_paradox} quantify this
discrepancy at high throughput ($>1$\,Gbps).

\begin{table}[htbp]
  \centering
  \caption{Specific Energy Consumption at high throughput ($>$\,1\,Gbps):
    median values.}
  \label{tab:sec}
  \begin{tabular}{lrr}
    \toprule
    \textbf{Metric} & \textbf{SD-Core DPDK} & \textbf{USR-UPF} \\
    \midrule
    Total power at $>$1\,Gbps (W)             & $0.824$ & $3.478$ \\
    \sectable{} (J/Gb) & $0.274$ & $1.615$ \\
    \secnet{} (J/Gb)   & $0.002$ & $1.607$ \\
    \bottomrule
  \end{tabular}
\end{table}

\begin{figure}[htbp]
  \centering
  \includegraphics[width=\linewidth]{\figpath{fig5_sec_paradox}}
  \caption{Specific Energy Consumption versus throughput.
    (a)~\sectable{}: DPDK appears efficient due to low power denominator
    relative to flat numerator.
    (b)~\secnet{}: DPDK\textquoteleft s \secnet{} collapses to near zero
    because its incremental power $P_{\mathrm{net}} \approx 0$.
    Both metrics understate the operational energy cost of running a
    poll-mode UPF. Extreme low-throughput outliers are clipped to the
    98th percentile for clarity.}
  \label{fig:sec_paradox}
\end{figure}

\subsection{Implications for Energy Metric Selection}

For a fair cross-variant comparison, \sectable{} is the more honest
metric: it reflects the full operational cost of sustaining throughput $\lambda$.
At $>1$\,Gbps, DPDK achieves 0.27\,J/Gb and USR achieves
1.61\,J/Gb---a 6$\times$ difference in total power
per bit.
However, total power consumption must also be considered in absolute terms:
at idle, DPDK draws 0.82\,W whereas USR draws less than
0.04\,W.
The preferable variant therefore depends on the traffic load distribution
in deployment.

% ── Section 6: Correlation Analysis ─────────────────────────────────────────
\section{Feature Correlation Analysis}
\label{sec:correlation}

\subsection{Correlates of Power Consumption}

Table~\ref{tab:correlations} lists the ten strongest Pearson correlations
with \texttt{power\_watts} in the pooled dataset (excluding features derived
directly from the target).

\begin{table}[htbp]
  \centering
  \caption{Top feature correlations with power consumption (pooled dataset).
    Rolling features are marked~$\dagger$ to indicate data-leakage risk when
    used as predictors.}
  \label{tab:correlations}
  \begin{tabular}{lcc}
    \toprule
    \textbf{Feature} & \textbf{Pearson $r$} & \textbf{Note} \\
    \midrule
    \texttt{power\_watts\_roll\_max\_30s}$^\dagger$  & $+0.998$ & Derived from target \\
    \texttt{cpu\_pct}                                & $+0.991$ & Strong legitimate predictor \\
    \texttt{power\_watts\_roll\_mean\_30s}$^\dagger$ & $+0.982$ & Derived from target \\
    \texttt{cpu\_per\_watt}                          & $-0.936$ & Derived from target \\
    \texttt{dl\_delay\_power\_product}               & $+0.797$ & Interaction feature \\
    \texttt{high\_delay\_frac}                       & $+0.768$ & QoS indicator \\
    \texttt{weighted\_mean\_delay\_us}               & $+0.735$ & Latency feature \\
    \texttt{ul\_delay\_power\_product}               & $+0.614$ & Interaction feature \\
    \texttt{l2l3\_overhead\_ratio}                   & $-0.561$ & Encapsulation efficiency \\
    \texttt{packets\_lost\_delta}                    & $+0.535$ & Loss indicator \\
    \bottomrule
  \end{tabular}
\end{table}

Figure~\ref{fig:correlation_heatmaps} shows the correlation structure
separately for each variant.
The two variants exhibit markedly different correlation patterns: in
\dpdk{}, power is weakly correlated with all traffic metrics ($r < 0.5$)
because the power signal is nearly flat; in \usr{}, the throughput--power
correlation is $r = 0.81$ and CPU--power is $r = 0.92$, making power highly
predictable from traffic observables.

\begin{figure}[htbp]
  \centering
  \includegraphics[width=\linewidth]{\figpath{fig6_correlation_heatmaps}}
  \caption{Pearson correlation heatmaps for a subset of key features,
    computed separately for each variant.
    The DPDK heatmap (left) shows uniformly low correlations with power,
    consistent with the near-constant power profile.
    The USR heatmap (right) reveals strong positive correlations between
    throughput, CPU, and power.}
  \label{fig:correlation_heatmaps}
\end{figure}

\subsection{Temporal Dynamics}

Figure~\ref{fig:temporal} shows representative power and throughput traces
within a single high-load run for each variant.
\dpdk{} exhibits essentially no power variation despite large throughput
swings; \usr{} tracks load closely until saturation, then power plateaus
while throughput (offered load) continues to increase.

\begin{figure}[htbp]
  \centering
  \includegraphics[width=\linewidth]{\figpath{fig4_temporal_traces}}
  \caption{Temporal power and throughput traces for representative
    high-load runs. SD-Core DPDK (top): power remains flat despite
    throughput variations. USR-UPF (bottom): power rises with load until
    the saturation plateau.}
  \label{fig:temporal}
\end{figure}

% ── Section 7: Implications ──────────────────────────────────────────────────
\section{Discussion and Implications}
\label{sec:implications}

\subsection{Architectural Trade-offs}

The measurements reveal a fundamental trade-off between the two I/O paradigms:

\begin{description}
  \item[Poll mode (DPDK)] delivers superior throughput ($5.06$ vs
    $2.33$\,Gbps), zero packet loss, and deterministically low latency, at
    the cost of a fixed power floor of 0.82\,W regardless of load.
    This is appropriate for high-throughput, latency-sensitive deployments
    where the UPF is expected to be consistently busy.
  \item[Interrupt mode (USR)] scales power proportionally to load,
    drawing near-zero power at idle, but saturates at
    $\approx$0.5\,Gbps and incurs severe delay and loss
    in the saturated regime.
    It is preferable only in low-to-medium load scenarios where
    energy-proportionality is prioritised over peak throughput.
\end{description}

\subsection{Safe Operating Regions}

Based on the per-run saturation analysis
(Section~\ref{sec:saturation}, Figure~\ref{fig:oai_per_run}),
the practical safe operating boundary for \usr{} is
$\lambda \leq 0.10\,Gbps$.
Above this threshold, packet loss becomes probable; above
0.57\,Gbps, it is near-certain.
No such constraint applies to \dpdk{} within the measurement range.

\subsection{Implications for Predictive Energy Modelling}

The near-constant power of \dpdk{} presents a challenge for regression-based
energy models: the target variable has a dynamic range of only 52\,mW,
approaching the noise floor of Scaphandre.
Models trained on DPDK data are expected to predict 0.82\,W for
almost any input; the information content of the traffic features is minimal.
Per-variant models are therefore strongly recommended over a single pooled model.

For \usr{}, the strong linear relationships ($r > 0.8$) between
throughput, CPU, and power suggest that simple linear or low-complexity
tree models may suffice in the unsaturated regime.
Training should be restricted to the safe operating region
($\lambda < 0.5$\,Gbps) unless the modelling objective explicitly includes
the saturated regime.

The rolling power features (\texttt{power\_watts\_roll\_mean\_30s},
\texttt{power\_watts\_roll\_max\_30s}) achieve very high correlations
with the power target ($r > 0.98$) and must be excluded as predictors
when the modelling goal is to predict power from non-power observables.

% ── Section 8: Summary ───────────────────────────────────────────────────────
\section{Summary}
\label{sec:summary}

This chapter characterised two 5G UPF implementations under controlled
load conditions.
The principal findings are:

\begin{enumerate}
  \item \textbf{DPDK is power-constant}: total power varies by only
    52\,mW over the full 0--5\,Gbps load range due to the
    CPU polling loop. This makes power prediction from traffic features
    statistically intractable.
  \item \textbf{USR is power-proportional but capacity-limited}: power
    scales linearly from $\approx 0$ to $\approx 3.4$\,W, but packet loss
    becomes near-universal above 0.57\,Gbps. The
    practical operating boundary is $\approx$0.1\,Gbps.
  \item \textbf{The SEC measurement paradox}: the \secnet{} metric
    misleadingly reports near-zero energy efficiency for DPDK because the
    idle power is almost equal to the load power. Total SEC (\sectable{})
    is the more operationally meaningful metric.
  \item \textbf{CPU is the strongest legitimate predictor of power}
    ($r = 0.99$ pooled), followed by delay-related features.
    Rolling power features are informative but constitute data leakage
    when used as predictors of power.
  \item \textbf{Separate per-variant models are necessary}: the two variants
    occupy distinct regimes with different predictive structures and cannot
    be modelled jointly without a variant indicator.
\end{enumerate}

% ── Bibliography placeholder ─────────────────────────────────────────────────
\bibliographystyle{IEEEtran}
% \bibliography{references}  % uncomment and point to your .bib file

\end{document}
